\chapter{About this book}

This mini-book is a print adaptation of the Wallet
Security framework from the \textbf{Security Frameworks},
an initiative by the \textbf{Security Alliance}. The web
pages were written for a site. This edition is rewritten
for a pocket book you can hand to someone who is about
to hold keys.

The source this edition follows is the reviewed personal
custody core: who holds keys, hot versus cold, hardware
setup, seed handling, signing verification, and
approvals. Software-wallet and hardware-wallet stub
pages, the signing-schemes stub, and the tools catalog
stay on the website. Empty chapters make a worse
companion than a short one.

Safe and Squads operator walkthroughs live in the
Protocol Multisig companion. This book will not pretend
to be that one.

This print cut changes order, cuts website navigation,
and turns the same controls into reading you can do on a
train. It does not invent new controls to fill gaps.

\section{Help us improve}

The live framework moves. If a control here is wrong,
incomplete, or now dangerous in your jurisdiction, fix it
upstream so the site and the next printing match.

{\ttfamily github.com/\\security-alliance/frameworks}

\keepblock{%
If you value this kind of resource, you can support the
work at our website.

{\centering\plainurl{securityalliance.org/donate}\par}
}

\keepblock{%
\section{Human safety comes first}

No asset, key, policy, or evidence goal is worth more
harm. If someone is being coerced, put the wrench down
and open the Physical Security companion. This book is
for the key. That book is for the person.
}

\chapter{Introduction}

Wallet security is control of keys and verification of
every signature. Hardware and multisig fail if you sign
attacker-controlled calldata because a website told you
to.

Cryptography protects keys from math. It does not
protect you from a fake dApp, a poisoned address, a
photo of a seed, or a hardware wallet that arrived
already set up.

This book is for people who will self-custody: beginners
with a hot wallet, operators moving savings onto
hardware, and anyone about to be the one person who can
spend. It is not a protocol treasury runbook.

Read it as a loop, not as a policy you recite once:

\begin{enumerate}
  \item Know who can spend, including the custodian you
  thought you did not have.
  \item Split spend from savings. Match value to
  exposure.
  \item Put the pile on hardware you bought and verified.
  \item Keep the seed offline. Treat digital exposure as
  full compromise.
  \item Never sign blindly. The device screen is the
  transaction.
  \item Cap approvals. Ignore mystery tokens.
  \item Add independent people before one key is
  life-changing.
\end{enumerate}

\section{Disclaimer}

This is not legal advice, and it is not a recovery
service. It will not match every chain, family, or
threat model. Some of the dollar figures, seed splits,
and slippage ranges in later chapters are examples from
the source, not defaults. Copying them blindly is how
you get a setup an attacker can predict.

A 3-piece seed split, a 25th word, or a custom
encryption scheme can lock you out harder than it locks
an attacker out. Design that kind of control with people
who do this work. Do not invent it at 2am.

The contents are free on the Security Frameworks site.
Suggest changes there.

\section{How this book is organized}

\begin{enumerate}
  \item \textbf{Who holds the keys.} Custody is who can
  spend.
  \item \textbf{Split the pile.} Hot for spend, cold for
  savings.
  \item \textbf{Hardware.} Buy direct, verify, sign on
  the device.
  \item \textbf{The seed.} Offline. Any leak is
  compromise.
  \item \textbf{Never sign blind.} Origin, address,
  function, amount.
  \item \textbf{Approvals.} Exact amounts, revoke,
  permits are not logins.
  \item \textbf{More than one key.} Independent signers,
  or theater.
  \item \textbf{New surfaces.} AA, 7702, TEEs. Power
  with extra failure modes.
  \item \textbf{If it breaks.} Move. Do not reuse.
\end{enumerate}

The last chapter is four short cards for the moment you
do not want to reread a section.
